\begin{workedexample}[Worked Example: Barkhausen criterion]
A feedback oscillator sustains oscillation when the loop gain magnitude
is at least unity and the loop phase is a multiple of $360^{\circ}$:
\[ |A\beta| \ge 1, \qquad \angle A\beta = 0^{\circ}\ (\mathrm{mod}\ 360^{\circ}). \]
At start-up the small-signal loop gain must exceed 1 so noise builds up; amplitude
then grows until nonlinearity trims the effective gain to exactly 1 in steady
state. The frequency is set where the loop phase is zero --- usually a high-$Q$
resonator.
\end{workedexample}

\begin{keytakeaways}
\begin{itemize}[leftmargin=1.4em,itemsep=2pt]
\item Oscillation needs loop gain $\ge 1$ and $0^{\circ}$ (mod $360^{\circ}$) loop phase (Barkhausen).
\item Start-up needs gain $>1$; steady state settles to unity via nonlinearity.
\item A higher-$Q$ tank gives better frequency stability and lower phase noise.
\end{itemize}
\end{keytakeaways}

\begin{exercises}
\item What loop phase shift sustains oscillation?
\item Why must the start-up loop gain exceed 1?
\item Does a higher-$Q$ tank help or hurt phase noise?
\end{exercises}

\furtherreading{Razavi~\cite{razavi} (ch.~8); Lee~\cite{lee} (ch.~14).}
